\section{Deterministic binary lifts}

Fix (n\geq 3). Write the visible state set as
\[
V=\mathbb Z/n\mathbb Z
\]
with successor
\[
s(i)=i+1.
\]
The lifted state set is
\[
\widetilde V=V\times\Ztwo.
\]
The nontrivial deck transformation is
\[
\deck(i,\epsilon)=(i,\epsilon+1).
\]

\begin{definition}
A deterministic binary lift of (s) is a map
\(
T:\widetilde V\to\widetilde V
\)
such that
\[
\pi T=s\pi,
\qquad
T\deck=\deck T,
\]
where \(\pi(i,\epsilon)=i\).
\end{definition}

The covering condition requires the first coordinate of (T(i,\epsilon)) to
be (i+1). Deck covariance then implies that the second coordinate differs
from (epsilon) by a quantity depending only on (i). Consequently there is
a unique voltage vector
\[
v=(v_0,\ldots,v_{n-1})\in(\Ztwo)^n
\]
such that
\[
T_v(i,\epsilon)=(i+1,\epsilon+v_i).
\]
Thus the (2^n) local voltage assignments describe all deterministic
binary lifts before gauge identification.

